\documentclass[11pt,a4paper]{article}
\usepackage[utf8]{inputenc}
\usepackage[T1]{fontenc}
\usepackage[catalan]{babel}
\usepackage{lmodern}
\usepackage[sfdefault,lining]{FiraSans}
\usepackage[a4paper,top=2.3cm,bottom=2.2cm,left=1.8cm,right=1.8cm,headheight=1.3cm,headsep=0.4cm,footskip=1cm]{geometry}
\usepackage{amsmath,amssymb}
\usepackage{icomma}
\usepackage{graphicx}
\usepackage[table]{xcolor}
\usepackage{tikz}
\usepackage{fontawesome5}
\usepackage{ccicons}
\usepackage{fancyhdr}
\usepackage{lastpage}
\usepackage{tasks}
\usepackage{needspace}
\usepackage[most]{tcolorbox}
\usepackage{array}
\usepackage[hidelinks,pdftitle={Solucions de la fitxa de la unitat 6: Proporcionalitat i percentatges. 2n d'ESO},pdfsubject={Matemàtiques 2n ESO},pdfkeywords={raons, proporcions, proporcionalitat directa, proporcionalitat inversa, regla de tres, proporcionalitat composta, repartiments proporcionals, percentatges, augments i disminucions percentuals}]{hyperref}

\definecolor{navy}{HTML}{1B2748}
\definecolor{accent}{HTML}{E8603A}
\definecolor{teal}{HTML}{1F9E90}
\definecolor{soft}{HTML}{EEF1F8}
\definecolor{muted}{HTML}{5D6685}

\setlength{\parindent}{0pt}
\setlength{\parskip}{3pt}
\settasks{label=\alph*), label-format=\bfseries\color{navy}, label-width=1.3em, item-indent=1.8em, column-sep=0.6em, after-item-skip=5pt, before-skip=3pt, after-skip=4pt}

% ---------- capçalera i peu ----------
\newcommand{\llicencia}{\href{https://creativecommons.org/licenses/by-nc-sa/4.0/deed.ca}{\ccbyncsa}\ \ Llicència Creative Commons Reconeixement-NoComercial-CompartirIgual 4.0 Internacional (CC BY-NC-SA 4.0)}
\pagestyle{fancy}
\fancyhf{}
\fancyhead[L]{\raisebox{-0.35\height}{\includegraphics[height=0.95cm]{logo-gva.png}}}
\fancyhead[C]{\small\color{muted}Solucions de la fitxa · Unitat 6: Proporcionalitat i percentatges · 2n d'ESO}
\fancyhead[R]{\raisebox{-0.35\height}{\includegraphics[height=0.95cm]{logo-vermellar.png}}}
\fancyfoot[L]{\scriptsize\color{muted}\llicencia}
\fancyfoot[R]{\small\color{muted}\thepage}
\renewcommand{\headrulewidth}{0.4pt}
\renewcommand{\headrule}{\color{muted!40}\hrule height \headrulewidth}
\fancypagestyle{primera}{\fancyhf{}\renewcommand{\headrulewidth}{0pt}
  \fancyfoot[L]{\scriptsize\color{muted}\llicencia}
  \fancyfoot[R]{\small\color{muted}\thepage}}

% ---------- nivells i exercicis ----------
\newcommand{\estrelles}[1]{{\color{accent}\ifcase#1\or\faStar\faStar[regular]\faStar[regular]\or\faStar\faStar\faStar[regular]\or\faStar\faStar\faStar\fi}}
\newcounter{exer}
\newcommand{\ex}[1]{\par\needspace{4\baselineskip}\medskip\stepcounter{exer}\noindent{\color{navy}\bfseries\theexer.}\ \ {\small\estrelles{#1}}\ \ }

\newcommand{\seccio}[2]{\par\needspace{10\baselineskip}\bigskip
  \begin{tcolorbox}[colback=navy,colframe=navy,coltext=white,arc=2mm,boxrule=0pt,left=3mm,right=3mm,top=1.2mm,bottom=1.2mm]
  {\large\bfseries #1\quad #2}
  \end{tcolorbox}\vspace{-2pt}}

\newcolumntype{C}[1]{>{\centering\arraybackslash}p{#1}}
\usepackage{pifont}
% quadre de recordatori (només a la fitxa d'exercicis)
\newcommand{\recorda}[1]{\par\smallskip\begin{tcolorbox}[colback=accent!7,colframe=accent!55,arc=2mm,boxrule=0.5pt,left=2.5mm,right=2.5mm,top=1mm,bottom=1mm]\small{\color{accent}\faLightbulb[regular]}\ \textbf{Recorda:}\ #1\end{tcolorbox}}
\newcommand{\res}[1]{\boxed{#1}}\newcommand{\rest}[1]{\fbox{#1}}

\def\fmneed{14}
% =====================================================================
% Capçalera FentMates de les fitxes: un color per a cada curs.
% Els colors són els de la pàgina inicial (LEVELS a assets/js/materials.js);
% el logo ve de _eines/logo/pdf/ (python3 _eines/logo/genera.py).
%
% Ús, al final del preàmbul de la unitat:
%   % =====================================================================
% Capçalera FentMates de les fitxes: un color per a cada curs.
% Els colors són els de la pàgina inicial (LEVELS a assets/js/materials.js);
% el logo ve de _eines/logo/pdf/ (python3 _eines/logo/genera.py).
%
% Ús, al final del preàmbul de la unitat:
%   % =====================================================================
% Capçalera FentMates de les fitxes: un color per a cada curs.
% Els colors són els de la pàgina inicial (LEVELS a assets/js/materials.js);
% el logo ve de _eines/logo/pdf/ (python3 _eines/logo/genera.py).
%
% Ús, al final del preàmbul de la unitat:
%   \input{../comu/fentmates.tex}
%   \fmcurs{1eso}
% i, al començament del document:
%   \capcalera{Matemàtiques}{1r ESO}{Fitxa d'exercicis · Unitat 2}{Divisibilitat}
% Opcional abans de \input: \def\fmneed{14} (espai mínim abans de cada secció).
% =====================================================================
\makeatletter
\def\fm@logodir{../../logo/pdf}
\providecommand{\fmneed}{10}
\def\fm@cursnom{}

% clau del curs, color i color del text damunt del color
\newcommand{\fm@curs}[3]{\@namedef{fm@color@#1}{#2}\@namedef{fm@text@#1}{#3}}
\fm@curs{1eso}{E6B84A}{navy}
\fm@curs{2eso}{3B82F6}{white}
\fm@curs{3eso}{FF7A50}{white}
\fm@curs{4eso-a}{A78BFA}{white}
\fm@curs{4eso-b}{6D28D9}{white}
\fm@curs{1batx-mi}{9B1C1C}{white}
\fm@curs{1batx-ccss}{EF6B6B}{white}
\fm@curs{2batx-mii}{1E6B47}{white}
\fm@curs{2batx-ccss}{2BB3A3}{white}

\newcommand{\fmcurs}[1]{%
  \@ifundefined{fm@color@#1}{\PackageError{fentmates}{Curs desconegut: #1}{}}{}%
  \edef\fm@tmp{\noexpand\definecolor{curs}{HTML}{\@nameuse{fm@color@#1}}%
               \noexpand\colorlet{curstext}{\@nameuse{fm@text@#1}}}\fm@tmp
  \colorlet{cursfosc}{curs!72!black}%
  \def\fm@logo{\fm@logodir/vertical-#1.pdf}%
  \def\fm@marca{\fm@logodir/marca-#1.pdf}}

% ---------- capçalera de la primera pàgina ----------
% #1 matèria (petit), #2 curs, #3 tipus de document i unitat, #4 títol
\newcommand{\capcalera}[4]{%
  \gdef\fm@cursnom{#2}%
  \par\noindent
  \begin{tcolorbox}[colback=curs,colframe=curs,boxrule=0pt,arc=2.2mm,width=3.5cm,height=2.15cm,
      valign=center,left=3mm,right=2mm,top=1mm,bottom=1mm,nobeforeafter,box align=center]
    \color{curstext}{\footnotesize #1}\par\vspace{1pt}{\fontsize{19}{21}\selectfont\bfseries #2\par}%
  \end{tcolorbox}\hspace{5mm}%
  \begin{minipage}[c]{\dimexpr\linewidth-3.5cm-5mm-2.9cm\relax}
    {\color{cursfosc}\small\bfseries\MakeUppercase{#3}}\par\vspace{2pt}
    {\color{navy}\fontsize{20}{23}\selectfont\bfseries #4\par}
  \end{minipage}\hfill
  \begin{minipage}[c]{2.6cm}\raggedleft\includegraphics[width=2.6cm]{\fm@logo}\end{minipage}\par
  \vspace{2pt}}

% ---------- barres de secció del color del curs ----------
\renewcommand{\seccio}[2]{\par\needspace{\fmneed\baselineskip}\bigskip
  \begin{tcolorbox}[colback=curs,colframe=curs,coltext=curstext,arc=2mm,boxrule=0pt,left=3mm,right=3mm,top=1.2mm,bottom=1.2mm]
  {\large\bfseries #1\quad #2}
  \end{tcolorbox}\vspace{-2pt}}

% ---------- estrelles de dificultat del color del curs (a Batxillerat no se'n posen) ----------
\@ifundefined{estrelles}{}{%
  \renewcommand{\estrelles}[1]{{\color{curs}\ifcase#1\or\faStar\faStar[regular]\faStar[regular]\or\faStar\faStar\faStar[regular]\or\faStar\faStar\faStar\fi}}}

% ---------- capçalera i peu de les pàgines ----------
% primera pàgina: logos institucionals al peu
\fancypagestyle{primera}{\fancyhf{}\renewcommand{\headrulewidth}{0pt}
  \fancyfoot[L]{\raisebox{-0.32\height}{\includegraphics[height=0.78cm]{logo-gva.png}}\hspace{5mm}%
    \raisebox{-0.32\height}{\includegraphics[height=0.78cm]{logo-vermellar.png}}}
  \fancyfoot[R]{\scriptsize\color{muted}\href{https://creativecommons.org/licenses/by-nc-sa/4.0/deed.ca}{\ccbyncsa}\ \ CC BY-NC-SA 4.0 · fentmates.org\qquad\small\thepage}}
% la resta: el símbol a l'esquerra i el curs a la dreta, amb el títol al centre com abans
\fancyhead[L]{\raisebox{-0.3\height}{\includegraphics[height=0.62cm]{\fm@marca}}}
\fancyhead[R]{\small\bfseries\color{cursfosc}\fm@cursnom}
\renewcommand{\headrule}{\color{curs!45}\hrule height \headrulewidth}
\makeatother

%   \fmcurs{1eso}
% i, al començament del document:
%   \capcalera{Matemàtiques}{1r ESO}{Fitxa d'exercicis · Unitat 2}{Divisibilitat}
% Opcional abans de \input: \def\fmneed{14} (espai mínim abans de cada secció).
% =====================================================================
\makeatletter
\def\fm@logodir{../../logo/pdf}
\providecommand{\fmneed}{10}
\def\fm@cursnom{}

% clau del curs, color i color del text damunt del color
\newcommand{\fm@curs}[3]{\@namedef{fm@color@#1}{#2}\@namedef{fm@text@#1}{#3}}
\fm@curs{1eso}{E6B84A}{navy}
\fm@curs{2eso}{3B82F6}{white}
\fm@curs{3eso}{FF7A50}{white}
\fm@curs{4eso-a}{A78BFA}{white}
\fm@curs{4eso-b}{6D28D9}{white}
\fm@curs{1batx-mi}{9B1C1C}{white}
\fm@curs{1batx-ccss}{EF6B6B}{white}
\fm@curs{2batx-mii}{1E6B47}{white}
\fm@curs{2batx-ccss}{2BB3A3}{white}

\newcommand{\fmcurs}[1]{%
  \@ifundefined{fm@color@#1}{\PackageError{fentmates}{Curs desconegut: #1}{}}{}%
  \edef\fm@tmp{\noexpand\definecolor{curs}{HTML}{\@nameuse{fm@color@#1}}%
               \noexpand\colorlet{curstext}{\@nameuse{fm@text@#1}}}\fm@tmp
  \colorlet{cursfosc}{curs!72!black}%
  \def\fm@logo{\fm@logodir/vertical-#1.pdf}%
  \def\fm@marca{\fm@logodir/marca-#1.pdf}}

% ---------- capçalera de la primera pàgina ----------
% #1 matèria (petit), #2 curs, #3 tipus de document i unitat, #4 títol
\newcommand{\capcalera}[4]{%
  \gdef\fm@cursnom{#2}%
  \par\noindent
  \begin{tcolorbox}[colback=curs,colframe=curs,boxrule=0pt,arc=2.2mm,width=3.5cm,height=2.15cm,
      valign=center,left=3mm,right=2mm,top=1mm,bottom=1mm,nobeforeafter,box align=center]
    \color{curstext}{\footnotesize #1}\par\vspace{1pt}{\fontsize{19}{21}\selectfont\bfseries #2\par}%
  \end{tcolorbox}\hspace{5mm}%
  \begin{minipage}[c]{\dimexpr\linewidth-3.5cm-5mm-2.9cm\relax}
    {\color{cursfosc}\small\bfseries\MakeUppercase{#3}}\par\vspace{2pt}
    {\color{navy}\fontsize{20}{23}\selectfont\bfseries #4\par}
  \end{minipage}\hfill
  \begin{minipage}[c]{2.6cm}\raggedleft\includegraphics[width=2.6cm]{\fm@logo}\end{minipage}\par
  \vspace{2pt}}

% ---------- barres de secció del color del curs ----------
\renewcommand{\seccio}[2]{\par\needspace{\fmneed\baselineskip}\bigskip
  \begin{tcolorbox}[colback=curs,colframe=curs,coltext=curstext,arc=2mm,boxrule=0pt,left=3mm,right=3mm,top=1.2mm,bottom=1.2mm]
  {\large\bfseries #1\quad #2}
  \end{tcolorbox}\vspace{-2pt}}

% ---------- estrelles de dificultat del color del curs (a Batxillerat no se'n posen) ----------
\@ifundefined{estrelles}{}{%
  \renewcommand{\estrelles}[1]{{\color{curs}\ifcase#1\or\faStar\faStar[regular]\faStar[regular]\or\faStar\faStar\faStar[regular]\or\faStar\faStar\faStar\fi}}}

% ---------- capçalera i peu de les pàgines ----------
% primera pàgina: logos institucionals al peu
\fancypagestyle{primera}{\fancyhf{}\renewcommand{\headrulewidth}{0pt}
  \fancyfoot[L]{\raisebox{-0.32\height}{\includegraphics[height=0.78cm]{logo-gva.png}}\hspace{5mm}%
    \raisebox{-0.32\height}{\includegraphics[height=0.78cm]{logo-vermellar.png}}}
  \fancyfoot[R]{\scriptsize\color{muted}\href{https://creativecommons.org/licenses/by-nc-sa/4.0/deed.ca}{\ccbyncsa}\ \ CC BY-NC-SA 4.0 · fentmates.org\qquad\small\thepage}}
% la resta: el símbol a l'esquerra i el curs a la dreta, amb el títol al centre com abans
\fancyhead[L]{\raisebox{-0.3\height}{\includegraphics[height=0.62cm]{\fm@marca}}}
\fancyhead[R]{\small\bfseries\color{cursfosc}\fm@cursnom}
\renewcommand{\headrule}{\color{curs!45}\hrule height \headrulewidth}
\makeatother

%   \fmcurs{1eso}
% i, al començament del document:
%   \capcalera{Matemàtiques}{1r ESO}{Fitxa d'exercicis · Unitat 2}{Divisibilitat}
% Opcional abans de \input: \def\fmneed{14} (espai mínim abans de cada secció).
% =====================================================================
\makeatletter
\def\fm@logodir{../../logo/pdf}
\providecommand{\fmneed}{10}
\def\fm@cursnom{}

% clau del curs, color i color del text damunt del color
\newcommand{\fm@curs}[3]{\@namedef{fm@color@#1}{#2}\@namedef{fm@text@#1}{#3}}
\fm@curs{1eso}{E6B84A}{navy}
\fm@curs{2eso}{3B82F6}{white}
\fm@curs{3eso}{FF7A50}{white}
\fm@curs{4eso-a}{A78BFA}{white}
\fm@curs{4eso-b}{6D28D9}{white}
\fm@curs{1batx-mi}{9B1C1C}{white}
\fm@curs{1batx-ccss}{EF6B6B}{white}
\fm@curs{2batx-mii}{1E6B47}{white}
\fm@curs{2batx-ccss}{2BB3A3}{white}

\newcommand{\fmcurs}[1]{%
  \@ifundefined{fm@color@#1}{\PackageError{fentmates}{Curs desconegut: #1}{}}{}%
  \edef\fm@tmp{\noexpand\definecolor{curs}{HTML}{\@nameuse{fm@color@#1}}%
               \noexpand\colorlet{curstext}{\@nameuse{fm@text@#1}}}\fm@tmp
  \colorlet{cursfosc}{curs!72!black}%
  \def\fm@logo{\fm@logodir/vertical-#1.pdf}%
  \def\fm@marca{\fm@logodir/marca-#1.pdf}}

% ---------- capçalera de la primera pàgina ----------
% #1 matèria (petit), #2 curs, #3 tipus de document i unitat, #4 títol
\newcommand{\capcalera}[4]{%
  \gdef\fm@cursnom{#2}%
  \par\noindent
  \begin{tcolorbox}[colback=curs,colframe=curs,boxrule=0pt,arc=2.2mm,width=3.5cm,height=2.15cm,
      valign=center,left=3mm,right=2mm,top=1mm,bottom=1mm,nobeforeafter,box align=center]
    \color{curstext}{\footnotesize #1}\par\vspace{1pt}{\fontsize{19}{21}\selectfont\bfseries #2\par}%
  \end{tcolorbox}\hspace{5mm}%
  \begin{minipage}[c]{\dimexpr\linewidth-3.5cm-5mm-2.9cm\relax}
    {\color{cursfosc}\small\bfseries\MakeUppercase{#3}}\par\vspace{2pt}
    {\color{navy}\fontsize{20}{23}\selectfont\bfseries #4\par}
  \end{minipage}\hfill
  \begin{minipage}[c]{2.6cm}\raggedleft\includegraphics[width=2.6cm]{\fm@logo}\end{minipage}\par
  \vspace{2pt}}

% ---------- barres de secció del color del curs ----------
\renewcommand{\seccio}[2]{\par\needspace{\fmneed\baselineskip}\bigskip
  \begin{tcolorbox}[colback=curs,colframe=curs,coltext=curstext,arc=2mm,boxrule=0pt,left=3mm,right=3mm,top=1.2mm,bottom=1.2mm]
  {\large\bfseries #1\quad #2}
  \end{tcolorbox}\vspace{-2pt}}

% ---------- estrelles de dificultat del color del curs (a Batxillerat no se'n posen) ----------
\@ifundefined{estrelles}{}{%
  \renewcommand{\estrelles}[1]{{\color{curs}\ifcase#1\or\faStar\faStar[regular]\faStar[regular]\or\faStar\faStar\faStar[regular]\or\faStar\faStar\faStar\fi}}}

% ---------- capçalera i peu de les pàgines ----------
% primera pàgina: logos institucionals al peu
\fancypagestyle{primera}{\fancyhf{}\renewcommand{\headrulewidth}{0pt}
  \fancyfoot[L]{\raisebox{-0.32\height}{\includegraphics[height=0.78cm]{logo-gva.png}}\hspace{5mm}%
    \raisebox{-0.32\height}{\includegraphics[height=0.78cm]{logo-vermellar.png}}}
  \fancyfoot[R]{\scriptsize\color{muted}\href{https://creativecommons.org/licenses/by-nc-sa/4.0/deed.ca}{\ccbyncsa}\ \ CC BY-NC-SA 4.0 · fentmates.org\qquad\small\thepage}}
% la resta: el símbol a l'esquerra i el curs a la dreta, amb el títol al centre com abans
\fancyhead[L]{\raisebox{-0.3\height}{\includegraphics[height=0.62cm]{\fm@marca}}}
\fancyhead[R]{\small\bfseries\color{cursfosc}\fm@cursnom}
\renewcommand{\headrule}{\color{curs!45}\hrule height \headrulewidth}
\makeatother

\fmcurs{2eso}

\begin{document}
\thispagestyle{primera}

\capcalera{Matemàtiques}{2n ESO}{Solucions de la fitxa · Unitat 6}{Proporcionalitat i percentatges}
\vspace{8pt}

\begin{tcolorbox}[colback=soft,colframe=soft,arc=2mm,left=3mm,right=3mm,top=1.5mm,bottom=1.5mm]
\small Intenta resoldre cada exercici abans de mirar-ne la solució. Si no t'ix el mateix resultat, repassa els passos o consulta els apunts interactius de la unitat: hi trobaràs calculadores de regla de tres, proporcionalitat composta, repartiments i percentatges que ho resolen pas a pas.
\end{tcolorbox}

% =====================================================
\seccio{1}{Raons i proporcions}

\ex{1} $\frac{360}{24}=\res{15}$: hi ha 15 alumnes per cada professor.

\ex{1} \begin{tasks}(2)
  \task $3\cdot20=60=4\cdot15$: \rest{sí}
  \task $5\cdot15=75\neq6\cdot12=72$: \rest{no}
  \task $2{,}5\cdot8=20=2\cdot10$: \rest{sí}
  \task $7\cdot28=196\neq9\cdot21=189$: \rest{no}
\end{tasks}

\ex{2} \begin{tasks}(4)
  \task $x=\frac{12\cdot15}{20}=\res{9}$
  \task $x=\frac{8\cdot9}{6}=\res{12}$
  \task $x=\frac{4{,}5\cdot8}{6}=\res{6}$
  \task $x=\frac{35\cdot6}{14}=\res{15}$
\end{tasks}

\ex{2} Xiques: $84-35=49$. Xics/xiques: $\frac{35}{49}=\res{\frac{5}{7}}$; xiques/total: $\frac{49}{84}=\res{\frac{7}{12}}$.

\ex{3} $\frac{5}{2}=\frac{750}{x}\Rightarrow x=\frac{2\cdot750}{5}=\res{300\ \text{mL}}$ de maduixes; en total, $\res{1050\ \text{mL}}$ de batut.

\ex{3} Per cada 3 anys de la menuda, la gran en té 4: $3+4=7$ parts; cada part són $28:7=4$ anys. \rest{12 i 16 anys} (i $\frac{12}{16}=\frac{3}{4}$ \ding{51}).


% =====================================================
\seccio{2}{Magnituds directament proporcionals}

\ex{1} \begin{tasks}(2)
  \task \rest{Sí} (preu = preu/kg · kg).
  \task \rest{No}.
  \task \rest{Sí} ($P=4\cdot c$).
  \task \rest{No} (amb 0 minuts ja paguem la quota).
\end{tasks}

\ex{1} \begin{tabular}{l|ccccc}Hores & 2 & 5 & 8 & \rest{12} & 20\\ \hline Km & 7 & \rest{17{,}5} & \rest{28} & 42 & \rest{70}\end{tabular}\quad Constant: $k=\frac{7}{2}=\res{3{,}5}$ km/h.

\ex{1} Una entrada: $30:4=7{,}50$ €. Set: $7{,}50\cdot7=\res{52{,}50\ \text{€}}$.

\ex{2} \begin{tasks}(2)
  \task $\begin{array}{rcl}100\ \text{km}&\longrightarrow&6{,}4\ \text{L}\\ 450\ \text{km}&\longrightarrow&x\end{array}$\quad $x=\frac{6{,}4\cdot450}{100}=\res{28{,}8\ \text{L}}$
  \task $\begin{array}{rcl}6{,}4\ \text{L}&\longrightarrow&100\ \text{km}\\ 40\ \text{L}&\longrightarrow&x\end{array}$\quad $x=\frac{100\cdot40}{6{,}4}=\res{625\ \text{km}}$
\end{tasks}

\ex{2} En 1 h: $1500:2{,}5=600$ botelles. En 8 h: $600\cdot8=\res{4800}$. Per a 2100: $2100:600=\res{3{,}5\ \text{h}}$ (3 h 30 min).

\ex{2} Per persona: 80 g d'arròs i 0,3 L de brou. Per a 10: \rest{800 g d'arròs i 3 L de brou}.

\ex{3} $1{,}6\cdot250=400$ cm $=4$ m; $\ 1{,}4\cdot250=350$ cm $=3{,}5$ m. Superfície: $4\cdot3{,}5=\res{14\ \text{m}^2}$.

\ex{3} a) 1a: 4, 8, 12, 24 €; \ 2a: 8, 11, 14, 23 €. b) \rest{Només en la primera} (quocient constant 4); en la segona $\frac{8}{1}\neq\frac{11}{2}$. c) Igualen a les 5 h (20 €); \rest{a partir de 6 hores} és més barata la segona.


% =====================================================
\seccio{3}{Magnituds inversament proporcionals}

\ex{1} \begin{tasks}(2)
  \task \rest{Inversa}
  \task \rest{Directa}
  \task \rest{Inversa}
  \task \rest{Cap de les dues} (augmenten juntes, però no exactament en la mateixa proporció)
\end{tasks}

\ex{1} \begin{tabular}{l|ccccc}Aixetes & 2 & 3 & 4 & \rest{6} & 12\\ \hline Hores & 18 & \rest{12} & \rest{9} & 6 & \rest{3}\end{tabular}\quad Constant: $k=2\cdot18=\res{36}$.

\ex{1} Una sola màquina: $12\cdot3=36$ h. Quatre: $36:4=\res{9\ \text{h}}$.

\ex{2} $\begin{array}{rcl}120\ \text{km/h}&\longrightarrow&3{,}5\ \text{h}\\ 140\ \text{km/h}&\longrightarrow&x\end{array}$\quad Inversa: $x=\frac{120\cdot3{,}5}{140}=\res{3\ \text{h}}$.

\ex{2} Constant: $40\cdot30=1200$. Amb 30 vaques: $1200:30=\res{40\ \text{dies}}$. Amb 48: $1200:48=\res{25\ \text{dies}}$.

\ex{2} El dipòsit té $45\cdot12=540$ L. $540:18=\res{30\ \text{cubs}}$; $\ 540:60=\res{9\ \text{L}}$.

\ex{3} Falten 12 dies de faena per a 10 obrers. Amb 8 obrers (inversa): $x=\frac{10\cdot12}{8}=15$ dies més. En total: $6+15=\res{21\ \text{dies}}$.

\ex{3} \begin{tasks}(1)
  \task Inversa: $\frac{30\cdot24}{20}=\res{36\ \text{files}}$.
  \task Directa: $\frac{5\cdot312}{120}=\res{13\ \text{caixes}}$.
  \task Inversa: $\frac{6\cdot85}{10}=\res{51\ \text{€}}$.
\end{tasks}


% =====================================================
\seccio{4}{Proporcionalitat composta}

\ex{1} D i D: $x=24\cdot\frac{9}{6}\cdot\frac{15}{10}=\res{54\ \text{kg}}$. (Un gos menja $24:6:10=0{,}4$ kg al dia.)

\ex{1} Una aixeta en una hora: $6000:4:5=300$ L. Llavors: $300\cdot6\cdot3=\res{5400\ \text{L}}$.

\ex{2} $\begin{array}{ccc}\text{Obrers} & \text{Hores/dia} & \text{Dies}\\ \hline 12 & 8 & 20\\ 16 & 6 & x\end{array}$\\[3pt]Si només canvien els obrers: $20\cdot\frac{12}{15}=\res{16\ \text{dies}}$. Amb 16 obrers i 6 h (I i I): $20\cdot\frac{12}{16}\cdot\frac{8}{6}=\res{20\ \text{dies}}$.

\ex{2} D i D: $540\cdot\frac{6}{4}\cdot\frac{20}{30}=\res{540\ \text{€}}$. (Casualment, el mateix!)

\ex{2} Més tones → més viatges (D); més camions → menys viatges cadascun (I): $8\cdot\frac{540}{360}\cdot\frac{5}{6}=\res{10\ \text{viatges}}$.

\ex{3} $\begin{array}{cccc}\text{Màquines} & \text{Hores/dia} & \text{Peces} & \text{Dies}\\ \hline 8 & 9 & 3600 & 5\\ 6 & 10 & 4800 & x\end{array}$\\[3pt]I, I i D: $x=5\cdot\frac{8}{6}\cdot\frac{9}{10}\cdot\frac{4800}{3600}=\res{8\ \text{dies}}$. (Una màquina fa $3600:8:9:5=10$ peces per hora.)

\ex{3} Més hectàrees → més tractors (D); menys dies → més tractors (I): $3\cdot\frac{30}{18}\cdot\frac{6}{5}=\res{6\ \text{tractors}}$.


% =====================================================
\seccio{5}{Repartiments proporcionals}

\ex{1} \begin{tasks}(2)
  \task $k=120$: \rest{360 i 480}
  \task $k=120$: \rest{240, 360 i 600}
\end{tasks}

\ex{1} $k=\frac{135}{45}=3$: \rest{60 € i 75 €}.

\ex{2} Simplifiquem: 2, 3 i 5. $k=\frac{7500}{10}=750$: \rest{1500 €, 2250 € i 3750 €}.

\ex{2} \begin{tasks}(2)
  \task $\frac{1}{2},\frac{1}{3}\to3,2$; $k=100$: \rest{300 i 200}
  \task $\frac{1}{2},\frac{1}{3},\frac{1}{4}\to6,4,3$; $k=100$: \rest{600, 400 i 300}
\end{tasks}

\ex{2} $k=\frac{1170}{26}=45$: \rest{270 €, 405 € i 495 €}. Invers: $\frac{1}{2},\frac{1}{3},\frac{1}{6}\to3,2,1$; $k=195$: \rest{585 €, 390 € i 195 €}.

\ex{3} $\frac{1}{1},\frac{1}{2},\frac{1}{3}\to6,3,2$; $k=\frac{2800}{11}$: \rest{1527,27 €, 763,64 € i 509,09 €} (arrodonit).

\ex{3} Entre el primer i el segon hi ha $5-3=2$ parts, que són 60 €: cada part val 30 €. Reben \rest{90 €, 150 € i 210 €}; en total, \rest{450 €}.


% =====================================================
\seccio{6}{Percentatges: concepte i càlcul}

\ex{1} {\renewcommand{\arraystretch}{2.3}\begin{tabular}[t]{|>{\centering\arraybackslash}m{3.9cm}|>{\centering\arraybackslash}m{3.9cm}|>{\centering\arraybackslash}m{3.9cm}|}\hline \textbf{Percentatge} & \textbf{Fracció irreductible} & \textbf{Decimal}\\ \hline\hline 30\,\% & $\res{\frac{3}{10}}$ & \rest{0,3}\\ \hline \rest{75\,\%} & $\dfrac{3}{4}$ & \rest{0,75}\\ \hline \rest{5\,\%} & $\res{\frac{1}{20}}$ & 0,05\\ \hline 2\,\% & $\res{\frac{1}{50}}$ & \rest{0,02}\\ \hline \rest{12{,}5\,\%} & $\dfrac{1}{8}$ & \rest{0,125}\\ \hline 150\,\% & $\res{\frac{3}{2}}$ & \rest{1,5}\\ \hline \end{tabular}}

\ex{1} \begin{tasks}(4)
  \task \res{48}
  \task \res{38}
  \task \res{90}
  \task \res{12}
  \task \res{9}
  \task \res{35}
  \task \res{9}
  \task \res{60}
\end{tasks}

\ex{1} \begin{tasks}(3)
  \task $350\cdot0{,}12=\res{42}$
  \task $84\cdot0{,}35=\res{29{,}4}$
  \task $1250\cdot0{,}04=\res{50}$
  \task $240\cdot0{,}075=\res{18}$
  \task $6000\cdot0{,}005=\res{30}$
  \task $65\cdot1{,}2=\res{78}$
\end{tasks}

\ex{2} $850\cdot0{,}46=\res{391}$ xics; $\ 850\cdot0{,}18=\res{153}$; $\ 850\cdot0{,}06=\res{51}$.

\ex{2} $58\cdot0{,}65=\res{37{,}7\ \text{kg}}$; $\ 72\cdot0{,}65=\res{46{,}8\ \text{kg}}$.

\ex{3} Dilluns: $2000\cdot0{,}3=600$; queden 1400. Dimarts: $1400\cdot0{,}25=350$; queden \rest{1050 L}. S'ha gastat $\frac{950}{2000}=\res{47{,}5\,\%}$ (no el 55 \%!).


% =====================================================
\seccio{7}{Quin percentatge representa? Càlcul del total}

\ex{1} \begin{tasks}(4)
  \task $\frac{12}{50}=\res{24\,\%}$
  \task $\frac{27}{36}=\res{75\,\%}$
  \task $\frac{45}{300}=\res{15\,\%}$
  \task $\frac{90}{72}=\res{125\,\%}$
\end{tasks}

\ex{1} \begin{tasks}(4)
  \task $36:0{,}2=\res{180}$
  \task $63:0{,}15=\res{420}$
  \task $240:0{,}75=\res{320}$
  \task $18:0{,}04=\res{450}$
\end{tasks}

\ex{2} $\frac{21}{28}=0{,}75$: \rest{75 \% aprovats} i \rest{25 \% suspesos}.

\ex{2} $182:0{,}35=\res{520\ \text{persones}}$.

\ex{2} Ella: $\frac{6}{15}=40\,\%$; ell: $\frac{7}{20}=35\,\%$. \rest{Ella}, tot i que ha encertat menys tirs.

\ex{3} Les 3 porcions són el $100-62{,}5=37{,}5\,\%$: total $=3:0{,}375=\res{8\ \text{porcions}}$. Cada porció: $\frac{1}{8}=\res{12{,}5\,\%}$.

\ex{3} $252:1{,}2=\res{210}$ bicicletes; aquest mes n'ha venut \rest{42 més}.


% =====================================================
\seccio{8}{Augments i disminucions percentuals}

\ex{1} \begin{tasks}(6)
  \task \rest{1,25}
  \task \rest{0,9}
  \task \rest{1,04}
  \task \rest{0,65}
  \task \rest{1,21}
  \task \rest{2}
\end{tasks}

\ex{1} \begin{tasks}(2)
  \task $90\cdot0{,}7=\res{63\ \text{€}}$
  \task $22\cdot1{,}04=\res{22{,}88\ \text{€}}$
  \task $1450\cdot1{,}03=\res{1493{,}50\ \text{€}}$
  \task $340\cdot0{,}85=\res{289\ \text{€}}$
\end{tasks}

\ex{2} $18\,500\cdot0{,}88=\res{16\,280\ \text{€}}$; ha baixat \rest{2220 €}.

\ex{2} \begin{tasks}(3)
  \task $68:0{,}8=\res{85\ \text{€}}$
  \task $157{,}30:1{,}21=\res{130\ \text{€}}$
  \task $3180:1{,}06=\res{3000}$
\end{tasks}

\ex{2} \begin{tasks}(3)
  \task $\frac{20}{25}=0{,}8$: \rest{$-20\,\%$}
  \task $\frac{43\,000}{40\,000}=1{,}075$: \rest{$+7{,}5\,\%$}
  \task $\frac{8}{6{,}4}=1{,}25$: \rest{$+25\,\%$}
\end{tasks}

\ex{3} $400\cdot1{,}1\cdot0{,}9=400\cdot0{,}99=\res{396\ \text{€}}$. Índex total $0{,}99$: ha baixat \rest{un 1 \%}.

\ex{3} a) No: $0{,}6\cdot0{,}9=0{,}54$, és un descompte del \rest{46 \%}. b) $37{,}80:0{,}54=\res{70\ \text{€}}$.

\ex{3} a) $250\cdot1{,}4\cdot0{,}8=250\cdot1{,}12=\res{280\ \text{€}}$. b) Índex $1{,}12$: guanya \rest{un 12 \%}.


\end{document}
